\color{unidarkgreen}{\subsection{Максимальная токовая защита}}
\color{black}

\begin{enumerate}[label=\arabic{section}.\arabic{subsection}.\arabic{enumi}, labelsep=4pt, leftmargin=0pt, itemindent=57pt, itemsep=0pt, parsep=0pt] % если что можно крутить itemsep


\item
Максимальная токовая защита (МТЗ) предназначена для резервирования действия основных 
быстродействующих защит при внутренних коротких замыканиях в БСК. Защита выполнена 
двухступенчатой. Измерительные цепи защиты включены на токи фаз А, В, С со стороны 
высоковольтных вводов БСК. При необходимости защита может быть включена на разность токов фаз 
(цифровая сборка токовых цепей в «треугольник»): 

\begin{equation}
\label{eq:Iabc}
\dot{I}'_{\textsl{A}} = \frac{ \dot{I}_{\textsl{A}} - \dot{I}_{\textsl{B}} }{\sqrt{3}} ; \dot{I}'_{\textsl{B}} = \frac{ \dot{I}_{\textsl{B}} - \dot{I}_{\textsl{C}} }{\sqrt{3}}; \dot{I}'_{\textsl{C}} = \frac{ \dot{I}_{\textsl{C}} - \dot{I}_{\textsl{A}}}{\sqrt{3}}
\end{equation}

\item
Защита имеет регулируемые значения по току и времени срабатывания. Каждая из ступеней 
защиты может выполнена с блокировкой по второй гармонике (отношению уровня второй гармоники к 
первой). Данная блокировка реализована на базе модуля \textbf{«ОВ БНТ ВВ»}. 


\begin{figure}[H]
\centering
\includegraphics[width=1\textwidth,height=1\textheight,keepaspectratio]{\fbpath/271.0281 МТЗ/img/MTZ_VV.pdf}
\captionof{figure}{Функциональный блок МТЗ}\label{pict:MTZ_VV}
\end{figure}


\small
\setlength{\extrarowheight}{0.05cm} % высота ячеек в таблице
\begin{longtable}
{|>{\centering\arraybackslash}m{4.5cm}|>{\centering\arraybackslash}m{4.5cm}|>{\centering\arraybackslash}m{3cm}|>{\centering\arraybackslash}m{3.5cm}|}
\caption{Перечень уставочных параметров МТЗ\hfill\vspace{-0.5\baselineskip}}\label{mtz:tbl1} \\ 
\hhline{|-|-|-|-|}
\rowcolor{gray!30}
Наименование уставки & Диапазон & Шаг &  Значение по умолчанию \\ 
\hhline{|-|-|-|-|}
\endfirsthead

\caption*{\hspace{3pt}\emph{Продолжение таблицы \ref{mtz:tbl1}\hfill\vspace{-0.5\baselineskip}}} \\ % сделано по ГОСТ 2.105 п.6.8.7
\hhline{|-|-|-|-|}
\rowcolor{gray!30}
Наименование уставки & Диапазон & Шаг &  Значение по умолчанию \\  
\hhline{|-|-|-|-|}
\endhead
\hhline{|-|-|-|-|}

\endfoot
\endlastfoot

%%%%%%%%%%%%содержимое%%%%%%%%%%%%%%

\end{longtable}
\setlength{\extrarowheight}{0.0cm} % возврат параметра высота ячеек в таблице
\normalsize

\end{enumerate} 